% !PRINTED 102
\weberpage{102}

We also say briefly that $\alpha$ is a root of $f(x)$. By § 4, $f(x)$ may then be divided by $x-\alpha$ without remainder, so that one may put
\begin{equation}
 f(x)=(x-\alpha)f_1(x),
\tag{2}
\end{equation}
where $f_1(x)$ is only of the $(n-1)$st degree and has the same first coefficient $a_0$ as $f(x)$, thus
\[
 f_1(x)=a_0x^{n-1}+a'_1x^{n-2}+\cdots .
\]
The coefficients of $f_1(x)$ are given in § 4, (6).

Every root of $f_1(x)$ is therefore at the same time a root of $f(x)$; and conversely every root of $f(x)$ is either $=\alpha$ or is a root of $f_1(x)$. If one root $\alpha$ of $f(x)$ is known, then the task of finding the remaining ones is reduced to the solution of an equation of the $(n-1)$st degree.

The aim of the considerations of this section consists in proving that every function $f(x)$ of the $n$th degree has at least one root. This theorem is called the \gesperrt{fundamental theorem of algebra}. First we draw from (2) the important consequence:

\medskip
\noindent I. \gesperrt{An equation of the $n$th degree cannot have more than $n$ roots.}
\medskip

For if $f(x)$ had more than $n$ roots, then $f_1(x)$, which is only of the $(n-1)$st degree, would have more than $n-1$ roots. But an equation of the first degree has not more than one root, whence the correctness of our theorem follows by complete induction.

Sometimes one also gives it the expression:

\medskip
\noindent II. \gesperrt{If a function $f(x)$ of the $n$th degree has more than $n$ roots, then all its coefficients must be zero.}
\medskip

If $\beta$ is a root of $f_1(x)$, then one may likewise put
\[
 f_1(x)=(x-\beta)f_2(x),\quad f(x)=(x-\alpha)(x-\beta)f_2(x),
\]
where $f_2(x)=a_0x^{n-2}+\cdots$ is only of the $(n-2)$nd degree. Thus, if one assumes that each of the functions obtained by this division, $f_1(x), f_2(x),\ldots$, has at least one root, one finally obtains
\begin{equation}
 f(x)=a_0(x-\alpha)(x-\beta)\cdots(x-\nu),
\tag{3}
\end{equation}
and we can therefore also express the theorem which we earlier described as the aim of our considerations as follows. It is to be proved:

\medskip
